% 04_limits.tex
\section{Limits of a change of variable}
\label{sec:limits}

The stereographic coordinate maps $[-1,1]$ onto the whole real line
and replaces the Chebyshev basis by the rational Chebyshev basis.
It is tempting to expect that a function which needs a high
polynomial degree in $\tilde x$ needs a low degree in $x$.  In
this section, we show that for functions of a block-encoded matrix
the degree is the same in both coordinates, and we identify the
setting in which the coordinate is of use.

We consider a block encoding $U$ of a Hermitian matrix $A/\lambda$
with eigenvalues $a_j\in[-1,1]$.  The QSVT circuit alternates $d$
applications of $U$ and $U^\dagger$ with phase rotations controlled
on the block-encoding ancillas, and it applies to the eigenvector
with eigenvalue $a_j$ a polynomial $P(a_j)$ of degree at most
$d$~\cite{GSLW2019,MartynRossiTanChuang2021}.  We recall the
argument in order to make one consequence explicit.  The
eigenvalue $a_j$ enters the circuit only through the rotation
angle $\theta_j=\arccos a_j$, and each query contributes one power
of $e^{\pm i\theta_j}$.  After $d$ queries every matrix element of
the circuit on the two-dimensional invariant subspace of $a_j$ is
therefore a trigonometric polynomial of degree at most $d$ in
$\theta_j$.  The block that the circuit encodes is a polynomial of
degree at most $d$ in $a_j=\cos\theta_j$, and the off-diagonal
elements are $\sin\theta_j=\sqrt{1-a_j^2}$ times a polynomial of
degree at most $d-1$.  Hence, the set of reachable transformations
is fixed by $d$ and is the same in every coordinate.

In particular, the expansion coefficients do not depend on the
coordinate.  Let $\theta=\arccos a$ and let
$\alpha=\cot\theta=a/\sqrt{1-a^2}$ be the stereographic coordinate
of Eq.~\eqref{eq:stereo-coord}.  Because
$T_k(a)=\TB_k(\alpha)=\cos k\theta$, every polynomial $P$ of degree
at most $d$ obeys
\begin{equation}
	P(a)=\sum_{k=0}^d c_k T_k(a)
	=\sum_{k=0}^d c_k\,\TB_k(\alpha)
	=\sum_{k=0}^d c_k\cos k\theta
	\label{eq:basis-indep}
\end{equation}
with one and the same set of coefficients $c_k$.  It follows that
the best approximation error of a target $f$ at degree $d$ is the
same in the Chebyshev, rational Chebyshev, and Fourier bases.

Since the error coincides at every degree, so does its rate of
decay with $d$.  For a target analytic near $[-1,1]$ this rate is
set by the Bernstein ellipse through the nearest singularity of
$f$~\cite[Ch.~8]{Trefethen2013}.  Under the stereographic map
these ellipses become the bipolar contours with foci $\pm i$ that
govern the convergence of rational Chebyshev
series~\cite{Boyd1990} (see also Sec.~17.7 of
Ref.~\cite{Boyd2001}), so that Boyd's convergence rate for the
rational Chebyshev series equals the Chebyshev rate.
Consequently, lower bounds for specific targets, such as those of
Sherstov~\cite{Sherstov2013} and the adversary bound of
Laneve~\cite{Laneve2025}, are unchanged by the change of variable.

The argument above takes the sample
points of the polynomial as given: they are the eigenvalues $a_j$
of the matrix $A$.  In state preparation, in contrast, there is no
given matrix.  The task is to write prescribed numbers $f(y_j)$
into the amplitudes of a register, and the points $y_j$ at which
$f$ is sampled are part of the design.  The argument does not
constrain this choice.  In Sec.~\ref{sec:stateprep}, we use this
freedom and choose the points such that $f$, written in the angle
variable, is a smooth function sampled on a uniform grid.
